\scanpage{170}
% archive sheet number: 169
\deleted{à \unsure{des cas} \illegible{1} \unsure{méthodes} \unsure{de descente} \unsure{fid.}}

D'ailleurs \deleted{\unsure{et}} comme une extension de $G$ par $\unsure{\gamma|\sigma}$ \unsure{n'a pas d'autom. non trivial}, on peut appliquer à $\operatorname{Ext}_{S\text{-gr}}(G,\gamma)$ les méthodes de descente [ainsi, pour $S'$ variable \inserted{\unsure{réd.}}, on obtient un foncteur en $S'$ compatible avec la descente fid. plate ....].

\unsure{Prenons}, p.ex. $G={\unsure{\mathbb G|G}}_a$, $\gamma=\mathbb{Z}/p\mathbb{Z}$, on trouve (\unsure{sous réserve|sans recours} \unsure{tt|et} \unsure{en|un} plus d'hypothèses noethériennes)
\[
\operatorname{Ext}_{S\text{-gr}}(\mathbb{G}_a,\mathbb{Z}/p\mathbb{Z})\simeq\Gamma(S,{\unsure{\mathscr O|\mathcal O|O}}_S^{p^{-\infty}}).
\]
% slanted
\margin{Pour la dém, cf exposé Groth. au séminaire \unsure{Chevalley|Chev.} sur les groupes algébriques I (Exposé 7)}
\[
\mathscr{O}_S^{p^{-\infty}}=\varinjlim_n(\mathscr{O}_S,n),
\]
avec l'hom $(\mathscr{O}_S,n)\longrightarrow(\mathscr{O}_S,m)$ défini par $x\longmapsto x^{p^{m-n}}$. On peut donc dire que le foncteur $\underline{\underline{\operatorname{Ext}}}_{S\text{-gr}}(\unsure{\mathbb{G}_S|\mathscr{O}_S},\mathbb{Z}/p\mathbb{Z})$ est représentable \deleted{par le} \inserted{dans la catégorie des} préschémas $p$-parfaits sur $S$.....
